\chapter*{前言}
\setchapterheader{前言}
\addcontentsline{toc}{chapter}{前言}
拿到一台机器人后，我们通常会先遇到三件事：认识它的部件，把它开起来，再让自己的程序与它通信。这本手册按这个顺序展开，先介绍硬件和日常操作，再说明开发环境，最后进入运动、音频和相机接口。

第一次使用 Bumi，可以从前两章开始，完成安装电池、连接 APP、起身和收纳。准备开发程序时，接着阅读第三至第五章，熟悉终端、Conda、网络和编译流程。调用出厂动作主要使用 HighController；部署自己的关节控制或运动策略，则使用 LowController。音视频接口和相机采集各自成章，便于按任务查阅。

命令需要在 Ubuntu 终端中执行。代码框中的路径和地址都对应具体用途：\code{~/noetix_sdk_bumi} 是本书示例使用的 SDK 根目录，\code{bumi_deploy} 是编译和运行共用的 Conda 环境，\code{192.168.55.102} 是运控板的默认地址。理解这些名字之后，再按自己的目录和网络填写命令，会更容易找到问题。

本文按\href{https://noetixrobotics.feishu.cn/docx/As2adPO27odPBfxQw1ucK6N4nNe}{《Bumi 综合使用文档》}重新编排，并结合其链接的编译环境、音视频说明及项目内 SDK 3.1.0 核对示例。正文中的表格用于查参数，命令框用于操作，提示框解释动作发生的条件。附录集中列出控制命令、工作模式、关节范围和视频入口。
